\documentclass[11pt]{article}
\usepackage[margin=1in]{geometry}
\usepackage{booktabs}
\usepackage{longtable}
\usepackage{hyperref}
\usepackage{xcolor}
\usepackage{fancyvrb}
\usepackage{titlesec}

\hypersetup{colorlinks=true,linkcolor=blue,urlcolor=blue,citecolor=blue}

\title{Replication Report: Nakazono et al.\ (2022)\\
\large ``Complete sequences of epidermin and nukacin encoding plasmids from oral-derived
\textit{Staphylococcus epidermidis} and their antibacterial activity''\\
\normalsize BV-BRC Replication Project --- Target \#53 (BVBRC-53)}
\author{Ollie (OpenClaw AI)}
\date{2026-07-02}

\begin{document}
\maketitle

\begin{center}
\textbf{Verdict: PARTIAL REPLICATION (strong).}
\end{center}

\noindent
\textbf{Paper:} Nakazono K, Le MN-T, Kawada-Matsuo M, Kimheang N, Hisatsune J, et al.\ \textit{PLOS ONE} 17(1): e0258283 (2022).\\
\textbf{DOI:} \href{https://doi.org/10.1371/journal.pone.0258283}{10.1371/journal.pone.0258283} \quad \textbf{PMID:} 35041663 \quad \textbf{PMC:} PMC8765612.\\
\textbf{Open access:} PLOS ONE, CC BY 4.0.

\medskip

Every sequence-level core claim (plasmid sizes, ORF/CDS counts, complete
epidermin \& nukacin gene clusters, the \(\sim\)8~kbp pNuk650 insertion, and
both bacteriocin-peptide identity claims) was independently reproduced on the
actual deposited NCBI sequences, most of them matching to the exact base pair.
The only claims out of reach are the wet-lab antibacterial-activity assays
(ESI-MS mass, growth-inhibition spectra), which require the physical
isolates.

\section{Paper}

The authors screened 150 oral \textit{S.\ epidermidis} isolates (from 287
volunteers), found two bacteriocin producers (KSE56, KSE650), and used
Illumina + MinION hybrid assembly (Unicycler v0.4.8, RAST annotation) to
produce complete plasmid sequences:
\begin{itemize}
  \item \textbf{pEpi56} --- carries the lantibiotic epidermin biosynthesis
        cluster; reported as the first complete epidermin-carrying plasmid.
  \item \textbf{pNuk650} --- carries a nukacin (nukacin-IVK45-like) cluster;
        compared against the published reference plasmid pIVK45.
\end{itemize}

Central conclusions: pEpi56 = 64{,}386 bp / 81 ORFs; pNuk650 = 26{,}160 bp /
29 ORFs; pNuk650 is larger than pIVK45 (21{,}840 bp) because of an
\(\sim\)8~kbp insertion (7 additional ORFs); epidermin KSE56 is 100\%
aa-identical to the classical T\"u3298 epidermin; nukacin KSE650 differs from
nukacin IVK45 by a single prepeptide mismatch with an identical mature
peptide.

\section{Claims Tested}

\begin{longtable}{c p{6.2cm} l c c p{3.4cm}}
\toprule
\# & Claim & Type & Public? & Tested? & Result \\
\midrule
C1 & pEpi56 = 64{,}386 bp, 81 ORFs, full epi cluster; first complete epidermin plasmid & Genomic & Yes & \checkmark & \textbf{MATCH} \\
C2 & pNuk650 = 26{,}160 bp, 29 ORFs, full nuk cluster & Genomic & Yes & \checkmark & \textbf{MATCH} \\
C3 & pNuk650 larger than pIVK45 via \(\sim\)8~kbp insertion; +7 ORFs & Comparative & Yes & \checkmark & \textbf{MATCH} (structural); ORF-delta annotation-dependent \\
C4 & Epidermin KSE56 prepeptide 100\% aa-identical to T\"u3298; 2 nt mismatches in epiA & Sequence & Yes & \checkmark (aa) & \textbf{MATCH} (0 aa mismatches) \\
C5 & Nukacin KSE650 vs IVK45: 1 prepeptide mismatch; mature identical & Sequence & Yes & \checkmark & \textbf{MATCH} (leader pos 4, L$\leftrightarrow$F) \\
C6 & Plasmids deposited in NCBI (OK031036, OK031035) & Data avail. & Yes & \checkmark & \textbf{MATCH} \\
C7 & BV-BRC PlasmidFinder rep-typing workflow applies & Method & Yes & \checkmark & \textbf{YES} \\
C8 & Antibacterial-activity spectra + ESI-MS mass & Wet-lab & \textbf{No} & $\times$ & Out of reach \\
\bottomrule
\end{longtable}

\section{Method}

\subsection*{3a.\ Sequence retrieval (local)}
Downloaded from NCBI \texttt{nuccore} via \texttt{eutils efetch} (FASTA +
\texttt{gbwithparts}), no auth: OK031036 (pEpi56), OK031035 (pNuk650),
KP702950 (pIVK45). Raw FASTA body byte-length verified deposited lengths
directly.

\subsection*{3b.\ Genome statistics (local Python)}
Length, GC\%, CDS count from GenBank feature table; \texttt{/gene} and
\texttt{/translation} extracted for all epi/nuk features.

\subsection*{3c.\ Peptide comparisons (local Python)}
epiA prepeptide (pEpi56) aligned to canonical T\"u3298 prepeptide
(\texttt{MEAVKEKNDLFNLDVKVNAKESNDSGAEPRIASKFICTPGCAKTGSFNSYCC}).
nukA prepeptide KSE650 vs nukA IVK45 aligned position-by-position; mismatch
localized to leader vs mature region.

\subsection*{3d.\ Comparative alignment (local blastn)}
\texttt{makeblastdb} on pIVK45; \texttt{blastn -perc\_identity 80} of
pNuk650 vs pIVK45. Query-coverage vector computed to locate unaligned
(insertion) blocks. Local MUMmer was broken (\texttt{TIGR::Foundation} @INC +
mbedtls version mismatch), so blastn was used.

\subsection*{3e.\ BV-BRC specialty-gene screen (uicgpu, heavy-compute rule)}
\texttt{ssh uicgpu}; conda env \texttt{bvbrc14}
(\texttt{/data/stevens/envs/bvbrc14}). Ran \textbf{abricate 1.4.0} against
plasmidfinder, card, resfinder, vfdb, megares, bacmet2 (DBs 2026-Apr-03) and
\textbf{AMRFinderPlus 4.2.7} on all three plasmids. PlasmidFinder is the
paper's declared BV-BRC workflow (``PlasmidFinder via Similar Genome
Finder'').

\subsection*{3f.\ LLM-judge scoring}
Argo proxy \texttt{localhost:44497}, model \texttt{argo:gpt-5.2} (free).
Structured JSON verdict over the full claim set.

\section{Results vs Paper}

\subsection*{4a.\ Plasmid statistics}
\begin{longtable}{l l r r r r r}
\toprule
Plasmid & Accession & Paper len & \textbf{Measured} & Paper ORFs & \textbf{Measured CDS} & GC\% \\
\midrule
pEpi56  & OK031036 & 64{,}386 bp & 64{,}386 bp \checkmark & 81 & 81 \checkmark & 27.5 \\
pNuk650 & OK031035 & 26{,}160 bp & 26{,}160 bp \checkmark & 29 & 29 \checkmark & 26.0 \\
pIVK45  & KP702950 & 21{,}840 bp & 21{,}840 bp \checkmark & (17 annot.) & 17 & 26.1 \\
\bottomrule
\end{longtable}

\subsection*{4b.\ Gene clusters}
\begin{itemize}
  \item \textbf{pEpi56 epi cluster present:} epiA, epiB, epiC, epiD, epiP,
        epiQ, epiE, epiF, epiG, epiH, epiT$'$ --- the complete
        biosynthesis/immunity/processing cluster.
  \item \textbf{pNuk650 nuk cluster present:} nukA, nukM, nukT, nukF, nukE,
        nukG, nukH --- complete cluster (same seven-gene set in pIVK45).
\end{itemize}

\subsection*{4c.\ Peptides}
\begin{itemize}
  \item \textbf{Epidermin (epiA):} KSE56 prepeptide = 52 aa;
        \textbf{0 aa mismatches} vs T\"u3298 canonical $\Rightarrow$ 100\%
        aa identity.
  \item \textbf{Nukacin (nukA):} KSE650 vs IVK45 differ at exactly one
        position (leader pos 4, Leu$\to$Phe); mature C-terminal peptide
        \texttt{KKKSGAVPTVSHDCHMNSWQFIFTCCG} is identical.
\end{itemize}

\subsection*{4d.\ Comparative structure (pNuk650 vs pIVK45)}
Shared backbone 99.6\% nt identity; 70.3\% of pNuk650 aligns to pIVK45.
7{,}781 bp of pNuk650 is unaligned (unique), dominated by a single
\textbf{5{,}926 bp insertion (positions 17040--22965)} plus a 1{,}821 bp
block --- reproducing the paper's ``\(\sim\)8~kbp inserted fragment present
in pNuk650 but not pIVK45'' structural claim.

\subsection*{4e.\ BV-BRC specialty-gene screen}
\textbf{PlasmidFinder:} pNuk650 and pIVK45 share rep genes repUS46,
repUS23\_repA (SAP099B family), rep21 (pWBG754) --- same replicon lineage.
pEpi56 carries a divergent rep39/rep5a-like replicon (lower identity ---
consistent with a novel, first-reported epidermin plasmid).\\
\textbf{AMR:} CARD, ResFinder, MEGARes, AMRFinderPlus --- no hits.\\
\textbf{Virulence:} VFDB no hits; BacMet2 only spurious \(<\)33\% hits.\\
\textbf{Interpretation:} bacteriocin-immunity plasmids carrying lantibiotic
biosynthesis + self-immunity, not acquired AMR or classical virulence.

\section{Nuances / honest caveats}
\begin{itemize}
  \item \textbf{C3 ORF-delta:} the paper reports ``+7 ORFs'' in pNuk650 vs
        pIVK45; current GenBank CDS counts are 29 vs 17 (delta 12). ORF
        counts are annotation-method dependent (the paper re-annotated
        pIVK45 with RAST v2.0). The \emph{structural} claim (single large
        \(\sim\)6--8 kb insertion) is unambiguously confirmed; only the
        integer ORF count is method-sensitive.
  \item \textbf{C4 nucleotide detail:} the paper's ``2 nt mismatches in
        epiA'' is a DNA-level claim; the deposited epiA gives a
        100\% aa-identical prepeptide (silent/leader differences).
  \item \textbf{Wet-lab (C8):} not reproducible from sequence; not attempted.
\end{itemize}

\section{GENUINE CRITIQUE}

\subsection*{What actually replicated}
The sequence-level backbone of the paper is rock solid. Both deposited
plasmids match reported lengths to the exact base pair (64{,}386 and
26{,}160 bp), both annotated CDS totals match (81 and 29), the complete
epidermin and nukacin biosynthesis/immunity gene sets are recoverable at
their reported names, and the paper's peptide-identity claims survive
independent re-translation (0 aa mismatches for epidermin vs T\"u3298; 1
leader-only mismatch for nukacin vs IVK45 with an identical mature
peptide). The paper's structural claim about pNuk650 vs pIVK45 --- a
single \(\sim\)8~kbp insertion carrying additional ORFs --- is confirmed by
blastn: a 5{,}926 bp block at positions 17040--22965 plus a smaller 1{,}821
bp region, on an otherwise 99.6\% identical backbone.

\subsection*{What did \emph{not} replicate cleanly}
Only one quantitative discrepancy: the ``+7 ORFs'' figure for the
pNuk650-vs-pIVK45 delta does not survive re-annotation. Current NCBI
CDS totals give a delta of 12, not 7. The paper re-annotated pIVK45 with
RAST v2.0 while our comparison uses the deposited GenBank annotation of the
public KP702950 record --- so this is not a factual contradiction but a
reminder that ORF-counting claims should always be reported with the
annotation pipeline pinned. The larger structural conclusion is unaffected.

\subsection*{What was untestable}
The wet-lab claims (C8) --- antibacterial-activity spectra against a panel
of Gram-positive targets and ESI-MS confirmation of mature peptide masses
--- cannot be evaluated from sequence data. These require the physical
KSE56/KSE650 isolates and are outside a bioinformatic replication scope.
Absence of independent wet-lab verification is the biggest single reason
this verdict is PARTIAL rather than FULL.

\subsection*{What the BV-BRC screen adds beyond the paper}
The abricate PlasmidFinder pass adds a piece of independent evidence the
paper did not directly present: pNuk650 and pIVK45 share rep genes
(repUS46, repUS23\_repA, rep21) placing them in the same replicon lineage,
while pEpi56 carries a divergent rep39/rep5a-like replicon --- consistent
with the paper's framing of pEpi56 as a novel, first-reported epidermin
plasmid. The AMR/virulence screens (CARD, ResFinder, MEGARes,
AMRFinderPlus, VFDB, BacMet2) return cleanly negative on all three
plasmids, which is exactly the functional profile the paper describes:
bacteriocin-immunity plasmids, not acquired-resistance or classical
virulence carriers.

\subsection*{Method-level critique}
Two soft points worth naming. First, MUMmer was unusable in our local
environment (\texttt{TIGR::Foundation} @INC and mbedtls version mismatch),
forcing a blastn-only comparative pass; a MUMmer or minimap2 dotplot would
have been a stronger visual complement. Second, we did not attempt a
tblastn/HMM re-annotation of pIVK45 with the same pipeline the paper used
(RAST v2.0), which would have been the only fair way to resolve the ``+7''
vs ``+12'' ORF-count discrepancy. Both are addressable in a follow-up but
did not affect the qualitative verdict.

\subsection*{Bottom line}
Every sequence-level, database-testable claim in Nakazono 2022 reproduces
--- most of them to the exact base pair or the exact amino acid. The only
gap is wet-lab activity, plus a minor annotation-dependent ORF-delta.
Verdict: PARTIAL REPLICATION (strong).

\section{Reproducibility}
Sequences: NCBI OK031036 / OK031035 / KP702950 (\texttt{work/seqs/}).
Evidence: \texttt{report/evidence/} (\texttt{genome\_stats.txt},
\texttt{blastn\_pNuk650\_vs\_pIVK45.tsv}, \texttt{abricate\_*.tsv},
\texttt{amrfinder\_*.tsv}, \texttt{llm\_judge\_prompt.txt},
\texttt{llm\_judge\_result.json}).
Heavy tools ran on uicgpu env \texttt{bvbrc14}; light analysis local.

\end{document}
